\makeabstract
{
Atomic and Molecular Ion Trap
}
{
Miles Antaya (1), David Hanneke (1)
}
{
\textsuperscript{1}Department of Physics and Astronomy, Amherst College, Amherst, MA, USA
}
{
We trap and image an atomic ion crystal mixed with oxygen inside a linear Paul trap to take precision measurements of the vibrational frequency of oxygen. To improve the trapping process, we are switching the coolant from Beryllium ions to Calcium ions to capture the oxygen at the center, instead of the outside of the crystal. However, this apparatus requires new lasers which may result in a loss of image resolution, making it harder to identify our trapped oxygen.
}
{
We positioned a 1951 USAF resolution test target beneath the aperture of the Andor iXon 885 Imager and reflected light from various LED and laser sources beneath the target. We were mainly concerned with comparing the 313 nm laser for exciting Beryllium (S\_1/2 → P\_1/2) to the 397 nm laser for exciting Calcium. We used a 395 nm LED instead for ease of this test, however we intend to use a 397 nm laser in the actual trap.
}
{
At both wavelengths, we are able to image stripes of width and distance apart 2.2µm from Group 7, Element 6 on the test target. Using a rectangular integrating brightness function across these stripes we are able to construct a brightness vs distance graph to identify the stripes by the peaks separated by troughs. While the range of gray value from trough to peak was around 3,000 to 12,000 on average for the 313 nm laser, the range for the 395 nm LED was 3,000 to 4,500 on average. Furthermore a qualitative inspection concluded that the test on the 395 nm LED appeared more indistinct, yet the lines were still distinguishable.
}
{
The lower range and lower average gray value for the 395 nm LED can likely be attributed to the less coherent light from the LED and a spherical aberration in the imaging system’s lenses, which result in a slightly blurry image. While we observed these measurement differences for stripes a few micrometers apart, identifying ions spaced less than 10 µm apart is still possible when using Calcium, with only a slight reduction in resolution. The difference is noticeable but will not significantly affect the functionality of our experiment.
}

\newpage
